% Converted from drafts/06_posthoc.md. Budget: 0.75 pages.
\section{Checks Added After Sealing}\label{sec:posthoc}

These checks were designed after the sealed results were known, to answer reviewers' questions.
None changes a sealed number; all are exploratory.

\textbf{What does hiding cost the cheater?} A5 gives up its saving on every request that looks like
a test. Replaying 172,418 real user requests from the WildChat
dataset through A5's classifier, only 3.3\% of requests (0.6\% of token cost) look like tests, so on
chat traffic A5 keeps over 99\% of the saving of always cheating. API traffic can differ: in Azure's
production code-completion trace~\cite{splitwise2024}, 76\% of requests generate at most 24 tokens, so
A5 could serve many honestly if they also repeat, which the trace does not record.

\textbf{Can disguised tests beat it?} When OTE asks its eight questions inside long, varied, everyday
requests that are never repeated, A5 no longer spots them, and OTE catches it in every session (AUROC
1.00, no honest session flagged), at 4.9$\times$ the cost per session. This holds by construction in
our simulator, which picks an answer from the question alone and replies with a bare token; a real
model answers in prose, and the surrounding request may shift its answers. IRIS finds real models
stay identifiable under disguise, but only partly across framings~\cite{iris2026}. Disguise also does
not help against a gateway that targets the auditor's account rather than its requests.

\textbf{Is the 10--15\% floor a property of auditing?} Varying the sealed choices, a power target of
50--90\% and a false-alarm budget of 1--10\% keep $\epsilon^*$ between 0.07 and 0.15 (up to 0.25
without IRIS-lite at a 1\% budget). The floor falls further only if sessions are pooled: fixing $k$ in advance and
running one one-sided Welch $t$-test of the $k$ sessions' mean score against the 300 honest
calibration sessions keeps the false-alarm rate at or below 5\% (checked on random splits of our 800
honest sessions). Pooled, GATEOPS catches $\epsilon = 0.02$--$0.03$ for under \$0.40 and BENCH
0.03, but pooled GATEOPS flags 100\% of honest load-balancing audits and pooled BENCH catches none
of the benchmark-aware A10. IRIS-lite reaches 0.03--0.07 but flags every honest change of sampling
settings (A4) and cannot separate the real models (Section~\ref{sec:real}). The only pooled auditor
that stays quiet on A4 and A11 and catches A10 is KBF, at $\epsilon = 0.07$ with 64 sessions (\$0.08,
plus \$0.40 for the 300 honest calibration sessions); A5 still beats it. Pooled timing also needs
its calibration taken at the same time as the audit: in the live run, honest latency drifted between
the first and second halves of 25 minutes (KS 0.20, $p = 0.046$).

\textbf{Does the confound depend on the invented timing?} Yes. The honest second provider behind A11
could not be measured, so its timing is hand-written. Set 18\% faster than the genuine model (the
measured gap between the two real models), GATEOPS still flags 74--100\% of load-balancing sessions;
set identical, it flags none. FUSE's dependence-safe average of e-values, which needs no
independence assumption, still flags 44\% of load-balancing sessions on the fitted pair and no
cheating session at $\epsilon = 0.10$.
